\chapter{Photons}

\section{What Is a Photon?}

Why can a sunburn come from invisible ultraviolet light but never from
the bright red glow of a heat lamp? Why do night-vision goggles work? The
answer is one of the strangest discoveries in physics: light comes in
lumps.

In the last chapters, we treated light as a wave, with each color a
different wavelength. But light also behaves like a stream of particles,
called \newterm{photons}. The term photons was coined by physicist Gilbert Newton Lewis, the same Lewis known for the Lewis dot structures we have talked about at the beginning of the sequence!

Every photon comes with an amount of energy that determines what color it is.
\index{photon}
When we are talking about light interacting with objects, your
intuition will be right more often if you think of light as a beam of
photons. Photons travel at the speed of light\index{speed of light}, $c \approx
2.998 \times 10^8$ m/s in a vacuum. Photons
\begin{itemize}
  \item behave both like particles and like waves;
  \item have no mass and no charge;
  \item are created when matter gives off light and destroyed when matter
        absorbs it.
\end{itemize}


\begin{mdframed}[style=important, frametitle={Formula for energy of a photon}]

For any given photon, its energy, $E$, is given by

$$E = hf = \frac{h c}{\lambda}$$

where $c$ is the speed of light, $\lambda$ is the wavelength of the
electromagnetic wave, and $h$ is Planck's constant: $6.63 \times 10^{-34}$ J$\cdot$s. (Recall
that $f = \frac{c}{\lambda}$.) Light energy comes in packets, or
\newterm{quanta}, of size $hf$, just as charge comes in packets of $e$. A
beam can carry 1 photon or 2 photons, but never 1.5.

For example, a red laser light has a wavelength of about 630 nm. So, the energy in each photon is:

$$\frac{(300 \times 10^6) (6.63 \times 10^{-34})}{630 \times 10^{-9}} \approx 3.2 \times 10^{-19} \text{ joules}$$

Joules are awkward for such tiny energies, so physicists use the
\newterm{electronvolt} (eV): the energy an electron gains crossing 1 volt,
$1.602 \times 10^{-19}$ J. In those units, $hc \approx 1240$ eV$\cdot$nm,
so $E = \frac{1240}{\lambda}$, with $\lambda$ in nanometers and $E$ in eV.
\end{mdframed}

\textbf{Worked example.} What is the energy of a photon of green light with
wavelength 500 nm?

\textbf{Solution.} $E = \frac{1240}{500} = 2.48$ eV.

Blue light has more energy per photon than red light; ultraviolet has more
still. That is why UV can damage the molecules in your skin and cause a
sunburn, while red light cannot, no matter how bright: each red photon is
too weak to break a molecular bond.

\subsubsection{The Photoelectric Effect}

Around 1900, physicists shined light on metal plates and measured the
electrons knocked loose. If light were just a wave, brighter light should
eventually shake electrons free, whatever its color, like waves slowly
pushing a boat to shore. But that is not what happened. Below a certain
frequency, no electrons came out, no matter how bright the light. Above
it, electrons came out instantly, even in dim light. This is the
\newterm{photoelectric effect}.\index{photoelectric effect}

In 1905, Albert Einstein explained it (and later won the Nobel Prize for
it). Light arrives in photons, and each electron is knocked loose by a
single photon. If one photon does not have enough energy, a million of them
will not help; brighter light only means more photons, not stronger ones.

The effect is governed by the conservation of energy. The energy of an
incoming photon ($E_{photon}$) must be at least the energy holding the
electron to the metal. This is described by $$E_{photon} = \Phi + K_{max}$$

where
\begin{itemize}
  \item $E_{\text{photon}} = hf$: The energy of the incoming photon, where $h$ is Planck's constant and $f$ is the frequency. 
  \item $\Phi$ (\newterm{Work Function})\index{work function}: The minimum energy required to eject an electron from a specific metal. Each material has a different ``cost'' to let an electron go.
  \item $f_0$ (\newterm{Threshold Frequency})\index{threshold frequency}: The specific frequency where the photon has exactly enough energy to match the work function ($\Phi = hf_0$). If $f < f_0$, no electrons are ejected. If $f = f_0$, electrons are released but have zero kinetic energy ($K_{\max} = 0$).
  \item $K_{\max}$: The ``leftover'' kinetic energy of the electron once released from the metal.
\end{itemize}
A closely related effect is how solar panels and digital camera sensors
turn light into electricity. In a solar cell, each photon of sunlight can
knock an electron loose inside the silicon. A built-in electric field in
the cell pushes the loose electrons in one direction, out through the
wires and around the circuit, creating a direct current. You will look at
camera sensors in depth in the camera chapter.

\begin{Exercise}[title={Ejected electrons}, label=ph_kmax]
Light with a wavelength of 400 nm shines on a metal with a work function
of 2.3 eV. What is the maximum kinetic energy of the ejected electrons?
\end{Exercise}
\begin{Answer}[ref=ph_kmax]
Each photon carries $\frac{1240}{400} = 3.10$ eV. After paying the 2.3 eV
work function, the electron keeps at most $3.10 - 2.3 = 0.80$ eV of
kinetic energy.
\end{Answer}
\subsubsection{Photons and Momentum}
In classical mechanics, we define momentum as the product of mass and velocity ($p = mv$). 

\index{photon!momentum of}This leads to a common question: if a photon has zero rest mass, how can it possess momentum? Energy and momentum are linked, even for massless particles. For a photon, the momentum ($p$) is inversely proportional to its wavelength ($\lambda$):
\begin{equation}
  p = \frac{h}{\lambda} \quad \text{or} \quad p = \frac{E}{c}
  \label{eq:photon_momentum}
\end{equation}
\textbf{Application:}
For example, \index{solar sail}solar sails are spacecraft propulsion systems in space. These large, mirror-like sails in space can intercept millions of photons. As these photons reflect off the surface, they undergo a change in momentum ($\Delta p$), which, according to Newton's Second Law ($F = \frac{\Delta p}{\Delta t}$), exerts a constant force on the spacecraft, allowing for fuel-free propulsion. 
\subsubsection{Photonic Interactions}
When a photon comes from the sun and hits an object, one of several
things can happen:

\begin{itemize}

  \item The energy of the photon is absorbed by the object. It makes the
        object a little warmer. If a large proportion of photons hitting the
        mass are absorbed like this, we say the object is ``black''.\index{absorption!photon}

  \item The photon bounces off the object. If the surface is very
        smooth, the photons bounce in a predictable manner, and we call
        this \newterm{reflection} and we say the object is ``shiny''.\index{reflection}

  \item If the surface is rough and the photons are not absorbed, the
        photons are scattered in random directions. We call this
        \newterm{diffuse reflection}. If most of the photons hitting an object are
        bounced in random directions, we say that the object is ``white''.\index{diffusion}

  \item The photon passes through the mass. If the mass has smooth
        surfaces and a consistent composition, the photons will pass through the
        mass in a predictable manner. We say that the mass is ``transparent''.\index{transparent}

  \item If the photons pass through, but in an unpredictable,
        scattering manner, we say the mass is ``translucent''. \index{translucent}

\end{itemize}

No object absorbs every photon, but chemists are always coming up with
``blacker'' materials. Vantablack, for example, is a super-black coating made of tiny carbon
tubes that absorbs 99.965\% of visible light. A crumpled sheet of foil
coated in it looks like a flat black hole.\index{Vantablack}

\begin{Exercise}[title={Photons and Color}, label=photon_color]

  There are red photons.

  Are there black photons?

  Are there white photons?

  Are there yellow photons?

\end{Exercise}
\begin{Answer}[ref=photon_color] \textit{Are there white photons?}  No. What we call ``white'' is a
  blend of photons that are several different colors.

  Some people like to say white light is the combination of all visible
  colors of photons in equal amounts. That seems oddly specific and unusual.

  Maybe it is better to imagine it from the human experience of white
  light. In our eyes, we have three different types of color-sensing
  cones, which generally correspond to the red, green, and
  blue regions of the spectrum. When all three are excited about equal
  amounts, humans experience that as white. On your computer screen,
  for example, what you see as white is just a blend of three colors of
  photons: a red, a green, and a blue.

  \textit{Are there black photons?}  No. What we call ``black'' is an
  absence of photons in the visible range.

  \textit{Are there yellow photons?} Yes! There is a region of the color
  spectrum that is yellow. It has a wavelength of about 580 nm. Photons
  at this energy level excite both our green-sensitive and red-sensitive
  cones (we will expand on this in the eye chapter!).
  Your computer monitor does not actually create light with a 580 nm
  wavelength. Instead, it creates red light and green light, which our
  eyes interpret as yellow.

\end{Answer}

\section{Exercises}

\begin{Exercise}[label=ph1, title={Energy of a photon}]
Find the energy, in electron volts, of a photon with wavelength 400 nm. Use $hc = 1240$ eV nm.
\end{Exercise}
\begin{Answer}[ref=ph1]
$E = \frac{hc}{\lambda} = \frac{1240 \text{ eV nm}}{400 \text{ nm}} = 3.10$ eV.
\end{Answer}

\begin{Exercise}[label=ph2, title={Momentum of a photon}]
Find the momentum of a photon with wavelength 630 nm. Use $h = 6.626 \times 10^{-34}$ J s.
\end{Exercise}
\begin{Answer}[ref=ph2]
$p = \frac{h}{\lambda} = \frac{6.626 \times 10^{-34}}{630 \times 10^{-9}} \approx 1.05 \times 10^{-27}$ kg m/s.
\end{Answer}

\begin{Exercise}[label=ph3, title={Threshold wavelength}]
A metal has a work function of 2.3 eV. What is the longest wavelength of light that can eject electrons from it? Use $hc = 1240$ eV nm.
\end{Exercise}
\begin{Answer}[ref=ph3]
The threshold wavelength is $\lambda_0 = \frac{hc}{\phi} = \frac{1240}{2.3} \approx 539$ nm. Light with a longer wavelength has too little energy per photon.
\end{Answer}

\begin{Exercise}[title={Red Light}, label=ph4]
What is the energy of a photon with wavelength $700$ nm?
\end{Exercise}
\begin{Answer}[ref=ph4]
$E = \frac{1240}{700} \approx 1.77$ eV.
\end{Answer}
